\section{GIỚI THIỆU ĐỀ TÀI}

\subsection{Bối cảnh và lý do chọn đề tài}

Trong mô hình doanh nghiệp nhiều chi nhánh, yêu cầu kết nối giữa các LAN không chỉ dừng ở việc ping thông mà còn phải đảm bảo độ ổn định, khả năng mở rộng và khả năng kiểm soát lưu lượng. Metro Ethernet là hướng triển khai phổ biến vì giữ được giao diện Ethernet quen thuộc ở phía khách hàng, đồng thời cho phép nhà cung cấp dịch vụ xây dựng mạng MAN có lõi tốc độ cao.

MPLS được sử dụng trong mạng nhà cung cấp để chuyển tiếp gói tin dựa trên nhãn thay vì chỉ tra cứu IP hop-by-hop. Cách tiếp cận này phù hợp với yêu cầu của đề tài: thiết kế backbone ISP có CE, PE, P router, kết nối ba chi nhánh doanh nghiệp và đánh giá hiệu năng khi mỗi chi nhánh dùng một kiến trúc LAN khác nhau.

\subsection{Mục tiêu nghiên cứu}

Project hướng đến bốn mục tiêu chính:
\begin{enumerate}
    \item Xây dựng topology Mininet gồm ba chi nhánh: Branch 1 Flat Network, Branch 2 Three-tier Network và Branch 3 Leaf-Spine Network.
    \item Triển khai provider backbone gồm CE, PE và P router, trong đó phần lõi dùng Linux kernel MPLS native để push, swap và pop label.
    \item Đo kiểm kết nối và hiệu năng bằng công cụ thật: ping, traceroute, iperf3 TCP, iperf3 UDP và tcpdump.
    \item Tự động hóa quy trình chạy test, sinh bảng kết quả, biểu đồ, GUI giám sát và báo cáo.
\end{enumerate}

\subsection{Phạm vi triển khai}

Mô hình được triển khai trên một máy Linux/Ubuntu có Mininet và Open vSwitch. Các router trong topology là Linux node có bật \texttt{net.ipv4.ip\_forward=1}. Mạng lõi MPLS dùng các module kernel \texttt{mpls\_router}, \texttt{mpls\_iptunnel}, FRRouting OSPF/LDP để hội tụ route và phân phối label động, cùng bảng \texttt{ip -f mpls route}. Ở biên dịch vụ, các CE được đưa vào service VPLS/L2VPN để lưu lượng liên chi nhánh đi qua data-plane L2VPN trên underlay MPLS.

\subsection{Đóng góp của project}

So với một topology Mininet cơ bản, project có các điểm mở rộng sau:
\begin{itemize}
    \item Có ba kiến trúc LAN khác nhau để so sánh ảnh hưởng đến hiệu năng.
    \item Có bằng chứng MPLS thật qua route table và tcpdump ethertype \texttt{0x8847}.
    \item Có UDP load sweep để đánh giá packet loss khi tải mạng tăng.
    \item Có GUI hỗ trợ demo từng phép đo theo cặp host được chọn.
    \item Có pipeline một lệnh để cleanup, vẽ topology, chạy benchmark, vẽ chart và tạo báo cáo.
\end{itemize}
